\subsection{Structure of fields with a non-ultrametric absolute value}

THEOREM 1 (Gelfand-Mazur). Let \(\mathbf{K}\) be an algebra over the field \(\mathbf{R}\) with the two following properties:

\begin{enumerate}
\def\labelenumi{(\arabic{enumi})}
\item
  K is a (not necessarily commutative) field.
\item
  There exists on K a norm \(x \mapsto \| x \|\) compatible with the algebra structure on K (General Topology, Chapter IX, § 3, no. 7, Definition 9).
\end{enumerate}

Then the algebra \(\mathbf{K}\) is isomorphic to one of the algebras \(\mathbf{R},\mathbf{C}\) or \(\mathbf{H}\) .

Recall (loc. cit.) that it may always be assumed that \(\|xy\|\leqslant\|x\|. \|y\|\) for all x,y in K. We shall give K the topology (compatible with the algebra structure) defined by the norm.

\textbf{(A) First case: \(\mathbf{K}\) is commutative and there exists \(j\in \mathbf{K}\) such that \(j^2 = -1\)}\par

Then there exists an isomorphism \(\sigma\) of the field \(\mathbf{C}\) onto a subfield of \(\mathbf{K}\) such that \(\sigma(\xi + i\eta) = \xi \cdot 1 + \eta \cdot j\) for \(\xi, \eta\) in \(\mathbf{R}\) . We shall prove by reductio ad absurdum that \(\mathbf{K} = \sigma(\mathbf{C})\) . Suppose then that there exists \(x \in \mathbf{K} - \sigma(\mathbf{C})\) ; for all \(z \in \mathbf{C}, x - \sigma(z)\) is therefore invertible in \(\mathbf{K}\) ; let us write \(\mathrm{F}(z) = (x - \sigma(z))^{-1}\) ; as \(\sigma\) is continuous and the inverse is continuous on \(\mathbf{K}\) (General Topology, Chapter IX, §3, no. 7, Proposition 13 applied to the completion algebra of \(\mathbf{K}\) ), \(\mathbf{F}\) is a continuous mapping of \(\mathbf{C}\) to \(\mathbf{K}\) . Moreover, we may write for \(z \neq 0\)

\[
\mathrm{F} (z) = (\sigma (z)) ^ {- 1} (x (\sigma (z)) ^ {- 1} - 1) ^ {- 1}.
\]

But, as \((\sigma(z))^{-1} = \sigma(z^{-1})\) tends to 0 as \(z\) tends to infinity in \(\mathbf{C}\) , it is seen that \(\mathbf{F}(z)\) tends to 0; in other words, \(z \mapsto \| \mathbf{F}(z) \|\) is a continuous real-valued function, \(\geqslant 0\) on \(\mathbf{C}\) , tending to 0 at the point at infinity and which can therefore be considered as a continuous function on the compact space \(\tilde{\mathbf{C}}\) obtained by adjoining to \(\mathbf{C}\) a point at infinity. The least upper bound \(\alpha\) of \(\| \mathbf{F} \|\) on \(\mathbf{C}\) is therefore finite and \(>0\) and the set \(\mathrm{P}\) of complex numbers \(z\) such that \(\| \mathrm{F}(z) \| = \alpha\) is closed and non-empty (General Topology, Chapter IV, §6, no. 1, Theorem 1).

Let \(z \in \mathbf{P}\) ; let us write \(y = x - \sigma(z)\) and let \(t\) be a complex number \(\neq 0\) such that \(\| \sigma(t) \| < \alpha^{-1}\) , whence \(\| \sigma(t).y^{-1} \| < 1\) by definition of \(\alpha\) . The sequence of the \((\sigma(t)y^{-1})^n\) and that of the \(n(\sigma(t)y^{-1})^n\) therefore tend to 0 in \(K\) as \(n\) tends to \(+\infty\) , for so do the corresponding sequences of norms in \(R\) . On the other hand note that for every polynomial \(\mathbf{H}(\mathbf{T}) = \prod_{k=1}^{p} (\mathbf{T} - \sigma(c_k))\) , where the \(c_k\) are distinct complex numbers, in the field \(\mathbf{K}(\mathbf{T})\) of rational functions

\[
\frac {\mathrm{H} ^ {\prime} (\mathrm{T})}{\mathrm{H} (\mathrm{T})} = \sum_ {k = 1} ^ {p} \frac {1}{\mathrm{T} - \sigma (c _ {k})}.\tag{5}
\]

We apply this formula to the polynomial

\[
\mathrm{H} (\mathrm{T}) = \mathrm{T} ^ {n} - (\sigma (t)) ^ {n} = \prod_ {k = 0} ^ {n - 1} \left(\mathrm{T} - \sigma \left(\omega_ {n} ^ {k} t\right)\right),
\]

where \(\omega_{n}=\exp(2\pi i/n)\) , and substitute for T the element \(y\in K\) , which is distinct from all the \(\sigma(\omega_{n}^{k}t)\) . It follows (in the commutative field K) that

\[
\frac {n y ^ {n - 1}}{y ^ {n} - (\sigma (t)) ^ {n}} = \frac {1}{y - \sigma (t)} + \sum_ {k = 1} ^ {n - 1} \frac {1}{y - \sigma \left(\omega_ {n} ^ {k} t\right)}.\tag{6}
\]

Taking account of the definitions of F and y, we obtain

\[
\begin{array}{r l} \mathrm{F} (z + t) + \sum_ {k = 1} ^ {n - 1} \mathrm{F} (z + \omega_ {n} ^ {k} t) - n \mathrm{F} (z) & \\ = \frac {n y ^ {n - 1}}{y ^ {n} - (\sigma (t)) ^ {n}} - \frac {n}{y} = \frac {1}{y} \cdot \frac {n (\sigma (t) y ^ {- 1}) ^ {n}}{1 - (\sigma (t) y ^ {- 1}) ^ {n}}. \end{array}\tag{7}
\]

But by virtue of the choice of t and the remarks made above, the last expression in (7) tends to 0 as n tends to \(+\infty\) ; hence

\[
\| \mathrm{F} (z + t) \| = \lim _ {n \to + \infty} \| n \mathrm{F} (z) - \sum_ {k = 1} ^ {n - 1} \mathrm{F} (z + \omega_ {n} ^ {k} t) \|.\tag{8}
\]

Now, \(\| \mathrm{F}(z)\| = \alpha\) and \(\| \mathrm{F}(z + \omega_n^kt)\| \leqslant \alpha\) by definition of \(\alpha\) , whence

\[
\left\| n \mathrm{F} (z) - \sum_ {k = 1} ^ {n - 1} \mathrm{F} (z + \omega_ {n} ^ {k} t) \right\| \geqslant
\]

\[
n \| \mathrm{F} (z) \| - \sum_ {k = 1} ^ {n - 1} \| \mathrm{F} (z + \omega_ {n} ^ {k} t) \| \geqslant n \alpha - (n - 1) \alpha = \alpha .
\]

Therefore by (8), letting \(n\) tend to \(+\infty\) , \(\| \mathbf{F}(z + t) \| \geqslant \alpha\) and by definition of \(\alpha\) this implies

\[
\| \mathrm{F} (z + t) \| = \alpha ,
\]

in other words \(z + t \in P\) . This proves that the set P is open in C; as it is also closed and non-empty and C is connected, P = C and \(\|F\|\) is therefore constant on C; as this function tends to 0 at the point at infinity, \(\|F(z)\| = 0\) in C and in particular \(\|F(0)\| = \|x^{-1}\| = 0\) , which is absurd.

\textbf{(B) Second case; \(\mathbf{K}\) is commutative and \(-1\) is not the square of an element of \(\mathbf{K}\)}\par

Let L be the commutative field obtained by adjoining to K a root j of \(T^{2} + 1\) ; L is a vector space over K admitting \((1, j)\) as a basis and L is obviously an algebra over R. Clearly the function \(x + yj \to \|x\| + \|y\|\) is a norm on L compatible with its structure as a vector space over R; on the other hand, for \(z = x + yj\) , \(z' = x' + y'j\) in L,

\[
\begin{array}{r l} \| z z ^ {\prime} \| & = \| x x ^ {\prime} - y y ^ {\prime} \| + \| x y ^ {\prime} + x ^ {\prime} y \| \leqslant \| x x ^ {\prime} \| + \| y y ^ {\prime} \| + \| x y ^ {\prime} \| + \| x ^ {\prime} y \| \\ & \leqslant \| x \| \cdot \| x ^ {\prime} \| + \| y \| \cdot \| y ^ {\prime} \| + \| x \| \cdot \| y ^ {\prime} \| + \| x ^ {\prime} \| \cdot \| y \| \\ & = (\| x \| + \| y \|) (\| x ^ {\prime} \| + \| y ^ {\prime} \|) = \| z \| \cdot \| z ^ {\prime} \|. \end{array}
\]

The norm thus defined is consequently compatible with the R-algebra structure on L. By case (A) L is an R-algebra isomorphic to C; now the only sub-R-algebra of C distinct from C is R and hence K is isomorphic to R.

\textbf{(C) Third case: K is not commutative}\par

Let Z be the centre of K and x an element of K not in Z; the subfield \(Z(x)\) of K is commutative and the norm induced by that on K is compatible with the R-algebra structure on \(Z(x)\) ; as \(Z \neq Z(x)\) and Z and \(Z(x)\) are R-algebras isomorphic to R or C by virtue of (A) and (B), Z is necessarily isomorphic to R and \(Z(x)\) to C. For all \(x \in K\) , \(Z(x)\) is therefore of rank \(\leqslant 2\) over Z. Now we have the following lemma:

LEMMA 1. Let D be a field with centre L such that, for all \(x \in \mathbf{D}\) , \(\mathbf{L}(x)\) is an extension of L of degree \(\leqslant m\) . Then the rank of D over L is \(\leqslant m^2\) .

We may obviously restrict our attention to the case where \(D \neq L\) . Then there exists in D a finite separable algebraic commutative extension E of L of degree >1 (Algebra, Chapter VIII, § 10, no. 3, Lemma 1); as \(E = L(x)\) for some suitable x in E (Algebra, Chapter V, § 7, no. 7, Proposition 12 and Chapter VII, § 5, no. 7), by hypothesis \([E: L] \leqslant m\) . Suppose the separable extension E is taken such that \([E: L]\) is finite and as great as possible and consider the centralizer \(E' \supset E\) of E in D, which is a field of centre E such that

\[
[ \mathbf {D} \colon \mathbf {E} ^ {\prime} ] = [ \mathbf {E} \colon \mathbf {L} ] \leqslant m
\]

(Algebra, Chapter VIII, § 10, no. 2, Theorem 2). If \(\mathbf{E} \neq \mathbf{E}'\) , there would exist in \(\mathbf{E}'\) a finite separable algebraic extension \(\mathbf{F}\) of \(\mathbf{E}\) of degree \(>1\) (Algebra, Chapter VIII, § 10, no. 3, Lemma 1); \(\mathbf{F}\) would therefore be a finite separable algebraic extension of \(\mathbf{L}\) (Algebra, Chapter V, § 7, no. 4, Proposition 7) of degree \(>[\mathrm{E:L}]\) , contrary to the definition of \(\mathbf{E}\) ; therefore \(\mathbf{E}' = \mathbf{E}\) , whence \([\mathrm{D:L}] = [\mathrm{D:E}][\mathrm{E:L}] \leqslant m^2\) .

Applying this lemma to K with m = 2, it is seen that K is a non-commutative extension field of R of finite rank and hence isomorphic to the field of quaternions H (Algebra, Chapter VIII, § 11, no. 2, Theorem 2).

Remark (1) We shall give in the chapter devoted to normed algebras a shorter proof of the Gelfand-Mazur Theorem which is valid for every Hausdorff locally convex topological algebra K over R and whose principle is the following: it is reduced (as in cases (B) and (C)) to the case where K is a commutative algebra over C; if \(x \in K - C.1\) , we consider as above the mapping \(z \mapsto (x - z.1)^{-1}\) of C to K, which is continuous and differentiable on C. For every element \(x'\) of the dual \(K'\) of the locally convex space K, \(z \mapsto \langle(x - z.1)^{-1}, x' \rangle\) is then a bounded integral function on C and therefore constant by Liouville's Theorem and we conclude as in part (A) of the proof of Theorem 1 that this necessarily implies \(\langle(x - z.1)^{-1}, x' \rangle = 0\) for all \(z \in C\) and all \(x' \in K'\) ; the Hahn-Banach Theorem shows that this conclusion is absurd, since \((x - z.1)^{-1} \neq 0\) . Note that the argument in part (A) of the proof of Theorem 1 differs from the above only in appearance, for this argument is only a special case of that which serves to prove the maximum principle for analytic functions, the summation over the roots of unity and passing to the limit being equivalent to calculating the integral \(\int_{\gamma} \frac{F(z + t) dt}{t}\) along a circle of centre 0 and the use of Cauchy's formula being avoided here, thanks to the particular form of the function F.

THEOREM 2 (Ostrowski). Let K be a (not necessarily commutative) field and f an element of \(\mathcal{V}(\mathbf{K})\) which is not an ultrametric absolute value. Then there exist a unique real number s > 0 and an isomorphism j of K onto an everywhere dense subfield of one of the fields R, C or H such that \(f(x) = |j(x)|^{s}\) for all \(x \in K (*)\) . For f to be an absolute value on K, it is necessary and sufficient that \(s \leqslant 1\) .

By no. 2, Corollary to Proposition 3, K is of characteristic 0 and hence an algebra over Q; for all \(x \in Q\) we write \(h(x) = f(x.1)\) ; clearly \(h \in \mathcal{V}(\mathbf{Q})\) and therefore Proposition 4 of no. 3 may be applied; neither of cases (i) and (ii) of the statement of this proposition can hold, for this would imply \(f(n.1) \leqslant 1\) for every integer n > 0 and f would be an ultrametric absolute value by virtue of no. 2, Proposition 3. Then there exists a real number \(s > 0\) such that \(h(x) = |x|^s\) for all \(x \in \mathbf{Q}\) , that is \(f(x,1) = |x|^s\) ; we write \(g = f^{1 / s}\) . Then \(g \in \mathcal{V}(\mathbf{K})\) and \(g(n,1) = n\) for every integer \(n\) ; Proposition 2 of no. 1 therefore shows that \(g\) is an absolute value on \(\mathbf{K}\) .

For \(x \in \mathbf{Q}\) and \(y \in \mathbf{K}\) , \(g(xy) = |x|g(y)\) and hence \(g\) is a norm on \(\mathbf{K}\) compatible with its \(\mathbf{Q}\) -algebra structure (with the usual absolute value on \(\mathbf{Q}\) ). The completion \(\hat{\mathbf{K}}\) of \(\mathbf{K}\) is therefore a normed algebra over \(\hat{\mathbf{Q}} = \mathbf{R}\) (General Topology, Chapter IX, § 3, no. 7); let \(\hat{g}\) be the norm on \(\hat{\mathbf{K}}\) the continuous extension of \(g\) . As \(g\) is an absolute value on \(\mathbf{K}\) , \(\hat{\mathbf{K}}\) is a field and \(\hat{g}\) an absolute value on \(\hat{\mathbf{K}}\) (General Topology, Chapter IX, § 3, no. 3, Proposition 6). By Theorem 1 there exists an \(\mathbf{R}\) -algebra isomorphism \(j\) of \(\hat{\mathbf{K}}\) onto one of the fields \(\mathbf{R}\) , \(\mathbf{C}\) or \(\mathbf{H}\) and \(g'(x) = |j(x)|\) is therefore an absolute value on \(\hat{\mathbf{K}}\) ; as \(\hat{\mathbf{K}}\) is finite-dimensional over \(\mathbf{R}\) and \(g'\) and \(\hat{g}\) coincide on the subfield \(\mathbf{R}\) . 1 of \(\hat{\mathbf{K}}\) , \(g' = \hat{g}\) by the following lemma:

LEMMA 2. Let L be a (not necessarily commutative) field and K a subfield of L such that L is a finite-dimensional left vector space over K. Let g be an absolute value on L and f its restriction to K. If K is complete and not discrete with respect to f, L is complete with respect to g; if further \(g'\) is another absolute value on L with the same restriction f to K, then \(g' = g\) .

As the topology defined by g is Hausdorff and compatible with the left vector K-space structure on L, the first assertion follows from Topological Vector Spaces, Chapter I, § 2, no. 3, Theorem 2. Moreover the topologies on L defined by g and \(g'\) are identical (loc. cit.); there therefore exists a real number s > 0 such that \(g' = g^{s}\) (General Topology, Chapter IX, § 3, no. 2, Proposition 5). Let x be an element of K such that \(f(x) \neq 1\) ; the equation \(g'(x) = g(x)\) proves that s = 1.

Returning to the proof of Theorem 2, it is seen that, if \(j\) denotes the restriction of \(\hat{j}\) to \(\mathbf{K}, j\) is an isomorphism of \(\mathbf{K}\) onto an everywhere dense subfield of \(\mathbf{R}\) , \(\mathbf{C}\) or \(\mathbf{H}\) and \(g(x) = |j(x)|\) for \(x \in \mathbf{K}\) , whence \(f(x) = |j(x)|^s\) .

Finally note that, if \(f\) is an absolute value on \(\mathbf{K}\) , \(h\) is an absolute value on \(\mathbf{Q}\) and \(s \leqslant 1\) by no. 3, Proposition 4; conversely, if \(s \leqslant 1\) , \(f = g^s\) is an absolute value on \(\mathbf{K}\) since \(g\) is (General Topology, Chapter IX, § 3, no. 2); this proves the last assertion of the statement.

\textbf{Remarks}\par

\begin{enumerate}
\def\labelenumi{(\arabic{enumi})}
\setcounter{enumi}{1}
\item
  If K is a field and a normed algebra over R, the norm is not necessarily an absolute value on K; for example, \(\xi + i\eta \to |\xi| + |\eta|\) is a norm on C compatible with its R-algebra structure.
\item
  For a proof of case (C) of Theorem 1 not using the general results of Algebra, Chapter VIII, see Exercise 2.
\end{enumerate}

